\section{Scope, Status, and Qualifications}

\subsection{Identity and coverage}

This is a source-first reference and usable bilingual prayer aid for the Angelus in the Latin Church's Catholic popular piety. It covers the received form, scriptural and doctrinal architecture, staged Western development, relation to the Regina Coeli, daily practice, and the current universal indulgence grant as printed in the Latin typical edition. It does not catalogue every territorial translation, bell pattern, religious-community custom, musical setting, private addition, or historical grant.

The governing devotional profile is the Marian-prayer addendum in \texttt{guidance/theology/mariology.md}. Liturgical claims are limited: the Angelus is treated as a pious exercise whose text and rhythm draw upon Scripture and liturgical prayer, not as an edition of the Roman Missal or Liturgy of the Hours. Current discipline is identified with the Apostolic Penitentiary's fourth \textit{Enchiridion Indulgentiarum}, promulgated in 1999 and checked through 21 July 2026. Territorial practice should be verified with the competent bishops' conference or local authority.

\subsection{Text, translation, and rights}

The reproduced Latin and English Angelus come from Dom Gaspar Lefebvre's \textit{Daily Missal with Vespers for Sundays and Feasts} (1925), printed p.~13; its abbreviation markers are retained and the full Latin and English Hail Mary on p.~5 of that same edition is printed separately below them. The scan pages were visually collated, and the checked file has SHA-256 \texttt{28c9046c5acfa19297ab9aeaa14ac4f8d8fcffab1307b51215b94fa4ca6172a2}. The title matter records historical ecclesiastical permissions for that edition. They do not approve this publication or make its English a current liturgical translation.

The 1925 edition is public domain in the United States by publication date. Its title page nevertheless says ``All rights reserved,'' and this study has not established authorship or non-US status for the exact historical English translation; non-US translation rights therefore remain unresolved. The received prayer wording is excluded from Triptych's CC BY 4.0 grant. Project-created exposition, arrangement, and design fall under that grant. No current proprietary vernacular liturgical translation is reproduced.

The current Apostolic Penitentiary and Holy See sources are cited and closely paraphrased rather than reproduced in English as a substitute edition. Scripture is cited by book and verse; no modern copyrighted biblical translation is reproduced in bulk.

\subsection{Historical and doctrinal limits}

The history presents a composite development. Medieval antecedents, a peace-bell lineage, a sixteenth-century developed form, later papal reception, and the modern prayer are not collapsed into one origin event. The study does not claim that one pope, friar, or council invented the complete Angelus.

Marian cooperation and intercession are interpreted under Christ's unique mediation. The study offers no apparition claim, promise of automatic spiritual effect, private timetable, or guarantee of peace or salvation. Its practical applications are editorial pastoral synthesis, not spiritual direction or binding law.

\subsection{Review state}

Completed internally: governing-source identification; prayer-text and page-image collation; claim-to-source audit; current-discipline and superseded-rule separation; rights review; generation-metadata validation; settled two-pass build; log, metadata, font, text, page-geometry, and PDF-structure checks; and visual inspection of every rendered page. Exact production details and the reviewed PDF hash are recorded in \texttt{research/scope.md}.

Outstanding: independent human collation of the Latin and historical English; independent Mariological, theological, liturgical, historical, Latin, legal, rights, and pastoral review; and any ecclesiastical review or approval. ``Source-audited'' describes internal evidence work only.
